% Análisis de División a/b bajo diferentes condiciones
% Los cuatro casos y sus combinaciones

\section{Análisis de División: $\frac{a}{b}$ bajo Diferentes Condiciones}

\subsection{Definición de los Cuatro Casos}

\begin{enumerate}
\item \textbf{Caso A:} $0 < a < 1$ (numerador es fracción positiva)
\item \textbf{Caso B:} $b > 1$ (denominador es mayor que 1)
\item \textbf{Caso C:} $a > 1$ (numerador es mayor que 1)
\item \textbf{Caso D:} $0 < b < 1$ (denominador es fracción positiva)
\end{enumerate}

\section*{Combinaciones de Casos y sus Efectos en $\frac{a}{b}$}

\subsection{Combinación 1: $0 < a < 1$ y $b > 1$}

\begin{caso}
Numerador fraccionario, denominador mayor que 1.
\end{caso}

\textbf{Efecto en $\frac{a}{b}$:} El resultado es \textbf{pequeño}.

\textbf{Justificación:}
\[
0 < a < 1 \implies \frac{a}{b} < \frac{1}{b}.
\]
Además, si $b > 1$, entonces $\dfrac{1}{b} < 1$, por lo tanto:
\[
\frac{a}{b} < 1.
\]
Más aún, cuanto mayor sea $b$, menor será $\dfrac{a}{b}$.

\begin{ejemplo}
$a = 0.5$, $b = 2$:
\[
\frac{a}{b} = \frac{0.5}{2} = 0.25.
\]
El resultado es más pequeño que ambos $a$ y $b$.
\end{ejemplo}

\begin{ejemplo}
$a = 0.3$, $b = 10$:
\[
\frac{a}{b} = \frac{0.3}{10} = 0.03.
\]
El resultado es muy pequeño.
\end{ejemplo}

\subsection{Combinación 2: $0 < a < 1$ y $0 < b < 1$}

\begin{caso}
Ambos, numerador y denominador, son fraccionarios.
\end{caso}

\textbf{Efecto en $\frac{a}{b}$:} El resultado \textbf{depende de la relación entre $a$ y $b$}.

\textbf{Justificación:}
\[
\frac{a}{b} = a \cdot \frac{1}{b}.
\]
Como $0 < b < 1$, se tiene $\dfrac{1}{b} > 1$. Por tanto, $\dfrac{a}{b}$ puede ser:
\begin{itemize}
\item Mayor que $a$ (si $b$ es muy pequeño)
\item Igual a $a$ (si $b = 1$, pero $b < 1$ por hipótesis)
\item Menor que $a$ nunca ocurre en esta combinación
\end{itemize}

En realidad: $\dfrac{a}{b} > a$ siempre que $0 < b < 1$.

\begin{ejemplo}
$a = 0.5$, $b = 0.2$:
\[
\frac{a}{b} = \frac{0.5}{0.2} = 2.5.
\]
El resultado es \textbf{mayor que} $a$ (amplificado por dividir por un número pequeño).
\end{ejemplo}

\begin{ejemplo}
$a = 0.3$, $b = 0.5$:
\[
\frac{a}{b} = \frac{0.3}{0.5} = 0.6.
\]
Nuevamente, $0.6 > 0.3 = a$.
\end{ejemplo}

\subsection{Combinación 3: $a > 1$ y $b > 1$}

\begin{caso}
Numerador y denominador ambos mayores que 1.
\end{caso}

\textbf{Efecto en $\frac{a}{b}$:} El resultado \textbf{depende de cuál es más grande}.

\textbf{Justificación:}
\[
\frac{a}{b} \begin{cases}
> 1 & \text{si } a > b, \\
= 1 & \text{si } a = b, \\
< 1 & \text{si } a < b.
\end{cases}
\]
Ambos son mayores que 1, pero el resultado puede estar en cualquier intervalo $(0, \infty)$ dependiendo de la magnitud relativa.

\begin{ejemplo}
$a = 5$, $b = 2$:
\[
\frac{a}{b} = \frac{5}{2} = 2.5 > 1.
\]
El numerador domina.
\end{ejemplo}

\begin{ejemplo}
$a = 2$, $b = 5$:
\[
\frac{a}{b} = \frac{2}{5} = 0.4 < 1.
\]
El denominador domina, aunque ambos son mayores que 1.
\end{ejemplo}

\subsection{Combinación 4: $a > 1$ y $0 < b < 1$}

\begin{caso}
Numerador mayor que 1, denominador fraccionario.
\end{caso}

\textbf{Efecto en $\frac{a}{b}$:} El resultado es \textbf{grande}.

\textbf{Justificación:}
\[
\frac{a}{b} = a \cdot \frac{1}{b}.
\]
Como $a > 1$ y $\dfrac{1}{b} > 1$ (pues $0 < b < 1$), obtenemos:
\[
\frac{a}{b} > 1 \cdot 1 = 1.
\]
De hecho, cuanto más pequeño sea $b$, mayor será $\dfrac{a}{b}$.

\begin{ejemplo}
$a = 4$, $b = 0.5$:
\[
\frac{a}{b} = \frac{4}{0.5} = 8.
\]
El resultado es grande por ambas razones: numerador amplifica, denominador pequeño amplifica aún más.
\end{ejemplo}

\begin{ejemplo}
$a = 3$, $b = 0.1$:
\[
\frac{a}{b} = \frac{3}{0.1} = 30.
\]
El resultado es muy grande.
\end{ejemplo}

\section{Resumen de Efectos: Tabla Comparativa}

\begin{center}
\begin{tabular}{|c|c|c|c|}
\hline
\textbf{Caso} & \textbf{Rango de $a$} & \textbf{Rango de $b$} & \textbf{Resultado $\frac{a}{b}$} \\
\hline
1 & $0 < a < 1$ & $b > 1$ & Muy pequeño: $\frac{a}{b} < a < 1$ \\
\hline
2 & $0 < a < 1$ & $0 < b < 1$ & Amplificado: $\frac{a}{b} > a$, puede ser $> 1$ \\
\hline
3 & $a > 1$ & $b > 1$ & Variable según $a$ vs $b$ \\
\hline
4 & $a > 1$ & $0 < b < 1$ & Muy grande: $\frac{a}{b} > a > 1$ \\
\hline
\end{tabular}
\end{center}

\section{Reglas Prácticas}

\begin{enumerate}
\item \textbf{Dividir por algo grande que 1 reduce el resultado.}
  \[
  \frac{a}{b} < a \quad \text{si } b > 1.
  \]

\item \textbf{Dividir por algo pequeño entre 0 y 1 amplifica el resultado.}
  \[
  \frac{a}{b} > a \quad \text{si } 0 < b < 1.
  \]

\item \textbf{Si el numerador es pequeño y el denominador es grande, el resultado es muy pequeño.}
  \[
  0 < a < 1, \ b > 1 \implies \frac{a}{b} \to 0.
  \]

\item \textbf{Si el numerador es grande y el denominador es pequeño, el resultado es muy grande.}
  \[
  a > 1, \ 0 < b < 1 \implies \frac{a}{b} \to \infty.
  \]

\item \textbf{Cuando ambos están en el mismo "lado" (ambos $>1$ o ambos $<1$), la magnitud relativa decide.}
  \[
  0 < a < 1, \ 0 < b < 1: \quad \frac{a}{b} > 1 \iff a > b.
  \]
  \[
  a > 1, \ b > 1: \quad \frac{a}{b} > 1 \iff a > b.
  \]
\end{enumerate}

\section{Visualización Conceptual}

\textbf{Escala de magnitudes:}

\vspace{0.5cm}

Muy pequeño \quad $\underbrace{0 < a < 1, b > 1}_{\text{Caso 1}}$ \quad $\underbrace{0 < a, b < 1}_{\text{Caso 2}}$ \quad $\underbrace{a, b > 1}_{\text{Caso 3}}$ \quad $\underbrace{a > 1, 0 < b < 1}_{\text{Caso 4}}$ \quad Muy grande

\vspace{0.5cm}

La posición en esta escala muestra cómo el resultado de $\dfrac{a}{b}$ varía dramáticamente según los rangos de $a$ y $b$.
